% !TEX program = xelatex
% 中文论文完整模板（参考文献内嵌，无需 .bib 文件）
% 编译方式：xelatex main.tex 跑两遍
\documentclass[12pt, a4paper]{ctexart}

% ==================== 页面设置 ====================
\usepackage[top=2.5cm, bottom=2.5cm, left=3cm, right=2.5cm]{geometry}
\usepackage{setspace}
\onehalfspacing

% ==================== 数学与符号 ====================
\usepackage{amsmath, amssymb, amsthm}
\usepackage{bm}

% ==================== 图表 ====================
\usepackage{graphicx}
\usepackage{booktabs}
\usepackage{multirow}
\usepackage{array}
\usepackage{caption}
\captionsetup{font=small, labelfont=bf}

% ==================== 超链接 ====================
\usepackage[colorlinks=true, linkcolor=black, citecolor=blue, urlcolor=blue]{hyperref}

% ==================== 定理环境 ====================
\newtheorem{theorem}{定理}[section]
\newtheorem{lemma}[theorem]{引理}
\newtheorem{definition}[theorem]{定义}
\newtheorem{remark}[theorem]{注}

% ==================== 自定义命令 ====================
\newcommand{\upcite}[1]{\textsuperscript{\cite{#1}}}

% ==================== 标题信息 ====================
\title{\heiti 论文题目：基于 XXX 的 XXX 方法研究}
\author{
    张三\thanks{作者简介：张三，硕士研究生，主要研究方向为 XXX，E-mail: xxx@xxx.com。} \\
    某某大学 某某学院，城市 邮编
}
\date{\today}

\begin{document}

\maketitle

% ==================== 摘要 ====================
\begin{abstract}
\noindent 这里是中文摘要。摘要应简要说明研究背景、目的、方法、主要结果和结论，一般 200--300 字。摘要中不应出现图表、公式和参考文献。

\vspace{1em}
\noindent\textbf{关键词：} 关键词1；关键词2；关键词3；关键词4

\vspace{1em}
\noindent\textbf{Abstract:} This is the English abstract. It should briefly describe the background, objective, methods, main results and conclusions of the paper.

\vspace{0.5em}
\noindent\textbf{Keywords:} keyword1; keyword2; keyword3; keyword4
\end{abstract}

\vspace{1em}

% ==================== 正文 ====================
\section{引言}
\label{sec:introduction}

近年来，XXX 问题受到了广泛关注\upcite{ref1}。然而，现有方法仍存在 XXX 不足\upcite{ref2}。本文针对该问题，提出了一种 XXX 方法。主要贡献如下：

\begin{enumerate}
    \item 提出了 XXX 模型；
    \item 设计了 XXX 算法；
    \item 在 XXX 数据集上验证了所提方法的有效性。
\end{enumerate}

\section{相关工作}
\label{sec:related}

\subsection{XXX 方法}
XXX 方法最早由 Smith 等人提出\upcite{ref2}。该方法通过 XXX 实现 XXX，但计算复杂度较高。

\subsection{XXX 方法}
针对上述问题，Li 等人提出了 XXX 改进方案\upcite{ref3}。

\section{问题描述与预备知识}
\label{sec:preliminaries}

\subsection{问题定义}
设数据集为 $\mathcal{D} = \{(\bm{x}_i, y_i)\}_{i=1}^{N}$，其中 $\bm{x}_i \in \mathbb{R}^d$ 为第 $i$ 个样本的特征向量，$y_i \in \{0,1\}$ 为对应标签。目标是学习一个映射 $f: \mathbb{R}^d \rightarrow \{0,1\}$，使得预测误差最小。

\subsection{符号说明}
本文使用的主要符号如表~\ref{tab:symbols} 所示。

\begin{table}[htbp]
\centering
\caption{符号说明}
\label{tab:symbols}
\begin{tabular}{cl}
\toprule
符号 & 含义 \\
\midrule
$\bm{x}_i$ & 第 $i$ 个样本的特征向量 \\
$y_i$ & 第 $i$ 个样本的标签 \\
$N$ & 样本总数 \\
$d$ & 特征维度 \\
$\mathcal{L}$ & 损失函数 \\
\bottomrule
\end{tabular}
\end{table}

\section{方法}
\label{sec:method}

\subsection{模型框架}
本文提出的模型框架如图~\ref{fig:framework} 所示。

\begin{figure}[htbp]
\centering
% \includegraphics[width=0.8\textwidth]{figures/framework.pdf}
\fbox{\parbox[c][5cm][c]{0.8\textwidth}{\centering 此处插入图片：figures/framework.pdf}}
\caption{模型框架图}
\label{fig:framework}
\end{figure}

\subsection{目标函数}
模型的目标函数定义为
\begin{equation}
\label{eq:objective}
\min_{\bm{\theta}} \mathcal{L}(\bm{\theta}) = \frac{1}{N} \sum_{i=1}^{N} \ell(f(\bm{x}_i; \bm{\theta}), y_i) + \lambda \|\bm{\theta}\|_2^2,
\end{equation}
其中，$\bm{\theta}$ 为模型参数，$\ell(\cdot)$ 为损失函数，$\lambda > 0$ 为正则化系数。

\subsection{优化算法}
目标函数~\eqref{eq:objective} 可通过梯度下降法求解：
\begin{equation}
\bm{\theta}_{t+1} = \bm{\theta}_t - \eta \nabla_{\bm{\theta}} \mathcal{L}(\bm{\theta}_t),
\end{equation}
其中 $\eta$ 为学习率。

\section{实验}
\label{sec:experiments}

\subsection{数据集}
实验采用 XXX 和 XXX 两个公开数据集，具体统计信息如表~\ref{tab:datasets} 所示。

\begin{table}[htbp]
\centering
\caption{数据集统计信息}
\label{tab:datasets}
\begin{tabular}{lccc}
\toprule
数据集 & 样本数 & 特征数 & 类别数 \\
\midrule
Dataset A & 1000 & 20 & 2 \\
Dataset B & 5000 & 50 & 5 \\
\bottomrule
\end{tabular}
\end{table}

\subsection{实验设置}
实验环境为 Python 3.9 + PyTorch 1.12，学习率设为 0.01，批量大小设为 64，迭代次数为 100。

\subsection{实验结果与分析}
不同方法在数据集上的实验结果如表~\ref{tab:results} 所示。

\begin{table}[htbp]
\centering
\caption{不同方法的性能对比}
\label{tab:results}
\begin{tabular}{lcc}
\toprule
方法 & 准确率（\%） & F1 值（\%） \\
\midrule
方法 A & 85.2 & 84.7 \\
方法 B & 87.6 & 86.9 \\
\textbf{本文方法} & \textbf{90.3} & \textbf{89.8} \\
\bottomrule
\end{tabular}
\end{table}

由表~\ref{tab:results} 可知，本文方法在准确率和 F1 值上均优于对比方法。

\section{结论}
\label{sec:conclusion}

本文针对 XXX 问题，提出了一种 XXX 方法。实验结果表明，所提方法在 XXX 数据集上取得了较好的性能。未来工作将围绕 XXX 展开进一步研究。

% ==================== 致谢 ====================
\section*{致谢}
感谢 XXX 对本研究的支持。

% ==================== 参考文献（内嵌） ====================
\begin{thebibliography}{99}

\bibitem{ref1}
Wang W, Li M. A survey on XXX[J]. Journal of XXX, 2023, 10(2): 1--20.

\bibitem{ref2}
Smith J, Brown A. XXX method for XXX[C]//Proceedings of the International Conference on XXX. 2022: 100--110.

\bibitem{ref3}
Li H. XXX theory and practice[M]. Beijing: XXX Press, 2021.

\end{thebibliography}

\end{document}